\textbf{Sparse Table}
Max/Min/GCD Query
\begin{lstlisting}
struct SparseTable {
  SparseTable(vector<int>& data)
  {
    int n = data.size();
    int maxLog = log2(n) + 1;
    table.assign(n, vector<int>(maxLog));
    for (int i = 0; i < n; ++i) {
      table[i][0] = data[i];
    }
    for (int j = 1; (1 << j) <= n; ++j) {
      for (int i = 0; (i + (1 << j) - 1) < n; ++i) {
        table[i][j] = max(table[i][j - 1], table[i +
        (1 << (j - 1))][j - 1]);
      }
    }
  }
  int query(int L, int R)
  {
    int j = log2(R - L + 1);
    return max(table[L][j], table[R - (1 << j) +
    1][j]);
  }
  vector<vector<int>> table;
};
\end{lstlisting}
\begin{lstlisting}
int main()
{
  vector<int> data = { 1, 3, 2, 7, 9, 11, 3, 5, 6, 8 };
  SparseTable st(data);
  cout << "Max in range [1, 4]: " << st.query(1, 4)
  << endl;
}
\end{lstlisting}

\textbf{Sparse Table 2D} (up to 500 x 500 is safe)
\begin{lstlisting}
struct SparseTable2D {
  static const int LG = 10;
  int n, m;
  vector<vector<int>> a;
  vector<int> lg2;
  vector<vector<vector<vector<int>>>> st;
  SparseTable2D(int n, int m) : n(n), m(m) {
    a.assign(n, vector<int>(m, 0));
    lg2.resize(max(n, m) + 1, 0);
    for (int i = 2; i <= max(n, m); i++) lg2[i] =
    lg2[i >> 1] + 1;
    st.assign(n, vector<vector<vector<int>>>(m,
    vector<vector<int>>(LG, vector<int>(LG,
    0))));
  }
  void set_val(int i, int j, int val) { a[i][j] = val; }
  void build() {
    for (int i = 0; i < n; i++)
      for (int j = 0; j < m; j++)
        st[i][j][0][0] = a[i][j];
    for (int a = 0; a < LG; a++) {
      for (int b = 0; b < LG; b++) {
        if (a + b == 0) continue;
        for (int i = 0; i + (1 << a) <= n; i++) {
          for (int j = 0; j + (1 << b) <= m; j++) {
            if (!a)
              st[i][j][a][b] = max(st[i][j][a][b-1],
              st[i][j + (1 << (b - 1))][a][b - 1]);
            else
              st[i][j][a][b] = max(st[i][j][a-1][b],
              st[i + (1 << (a - 1))][j][a - 1][b]);
          }
        }
      }
    }
  }
  int query(int x1, int y1, int x2, int y2) {
    int a = lg2[x2 - x1 + 1], b = lg2[y2 - y1 + 1];
    int m1 = max(st[x1][y1][a][b], st[x2 - (1 << a)
    + 1][y1][a][b]);
    int m2 = max(st[x1][y2 - (1 << b) + 1][a][b],
    st[x2 - (1 << a) + 1][y2 - (1 << b) + 1][a][b]);
    return max(m1, m2);
  }
};
\end{lstlisting}
\begin{lstlisting}
// st.set_val(i, j, grid[i][j]);
// st.build(); // st.query(0, 0, 1, 2); // max in
// submatrix [(0,0) to (1,2)]
\end{lstlisting}

\switchcolumn

\textbf{Merge Sort Tree}
\begin{lstlisting}
struct MergeSortTree {
  vector<vector<int>> tree;
  vector<int> arr;
  int n;
  MergeSortTree(vector<int>& a)
  {
    arr = a, n = a.size();
    tree.resize(4 * n);
    build(1, 0, n - 1);
  }
  void build(int node, int l, int r)
  {
    if (l == r) {
      tree[node].push_back(arr[l]);
      return;
    }
    int mid = (l + r) / 2;
    build(node * 2, l, mid), build(node * 2 + 1,
    mid + 1, r);
    merge(tree[node * 2].begin(), tree[node *
    2].end(), tree[node * 2 + 1].begin(), tree[node * 2
    + 1].end(), back_inserter(tree[node]));
  }
  int query(int node, int l, int r, int ql, int qr, int
  k)
  {
    if (l > qr || r < ql)
      return 0;
    if (l >= ql && r <= qr) {
      int hi = tree[node].size() - 1, lo = 0, ind =
      -1;
      // while (hi >= lo) { } // Modify your
      // function here
      return ind;
    }
    int mid = (l + r) / 2;
    return query(node * 2, l, mid, ql, qr, k) +
    query(node * 2 + 1, mid + 1, r, ql, qr, k);
  }
};
\end{lstlisting}

Find the number of distinct elements in a range.
\begin{lstlisting}
void merge(vector<int> tree[], int treeNode)
{
  int len1 = tree[2 * treeNode].size();
  int len2 = tree[2 * treeNode + 1].size();
  int index1 = 0, index2 = 0;
  while (index1 < len1 && index2 < len2) {
    if (tree[2 * treeNode][index1]
    > tree[2 * treeNode + 1][index2]) {
      tree[treeNode].push_back(
      tree[2 * treeNode + 1][index2]);
      index2++;
    }
    else {
      tree[treeNode].push_back(
      tree[2 * treeNode][index1]);
      index1++;
    }
  }
  while (index1 < len1) {
    tree[treeNode].push_back(
    tree[2 * treeNode][index1]);
    index1++;
  }
  while (index2 < len2) {
    tree[treeNode].push_back(
    tree[2 * treeNode + 1][index2]);
    index2++;
  }
  return;
}
\end{lstlisting}

\begin{lstlisting}
void build(vector<int> tree[], int* arr, int start,
int end,
int treeNode)
{
  if (start == end) {
    tree[treeNode].push_back(arr[start]);
    return;
  }
  int mid = (start + end) / 2;
  build(tree, arr, start, mid, 2 * treeNode);
  build(tree, arr, mid + 1, end, 2 * treeNode + 1);
  merge(tree, treeNode);
  return;
}
\end{lstlisting}

\switchcolumn

\begin{lstlisting}
int query(vector<int> tree[], int treeNode, int
start,
int end, int left, int right)
{
  if (start > right || end < left) {
    return 0;
  }
  if (start >= left && end <= right) {
    return tree[treeNode].end()
    - upper_bound(tree[treeNode].begin(),
    tree[treeNode].end(), right);
  }
  int mid = (start + end) / 2;
  int op1 = query(tree, 2 * treeNode, start, mid,
  left,
  right);
  int op2 = query(tree, 2 * treeNode + 1, mid + 1,
  end,
  left, right);
  return op1 + op2;
}
int main()
{
  int n; cin >> n;
  int arr[n];
  for(int i=0;i<n;i++) cin >> arr[i];
  int next_right[n];
  vector<int> tree[4 * n];
  unordered_map<int, int> ump;
  for (int i = n - 1; i >= 0; i--) {
    if (ump[arr[i]] == 0) {
      next_right[i] = n;
      ump[arr[i]] = i;
    }
    else {
      next_right[i] = ump[arr[i]];
      ump[arr[i]] = i;
    }
  }
  build(tree, next_right, 0, n - 1, 1);
  int q; cin >> q;
  while(q--){
    int l, r; cin >> l >> r;
    l-- , r--;
    cout << query(tree,1,0,n-1,l,r) << "\n";
  }
  return 0;
}
\end{lstlisting}

\textbf{Kth smallest value in a range (SegTree)}
\begin{lstlisting}
#define N (int)1e5
vector<int> seg[N];
void build(int ci, int st, int end,
pair<int, int>* B)
{
  if (st == end) {
    seg[ci].push_back(B[st].second);
    return;
  }
  int mid = (st + end) / 2;
  build(2 * ci + 1, st, mid, B);
  build(2 * ci + 2, mid + 1, end, B);
  merge(seg[2 * ci + 1].begin(), seg[2 * ci +
  1].end(),
  seg[2 * ci + 2].begin(), seg[2 * ci + 2].end(),
  back_inserter(seg[ci]));
}
\end{lstlisting}

